\section{Validación y reproducibilidad}
\label{sec:reproducibilidad}

Los resultados se acompañan de una cadena de validación que incluye una
implementación funcional independiente, pruebas automatizadas, testigos
reconstruibles, archivos estructurados y versiones identificables mediante
Git.

El propósito de esta sección es describir los elementos necesarios para
verificar las afirmaciones principales del trabajo sin repetir los detalles
de arquitectura presentados en la
Sección~\ref{sec:arquitectura_implementacion}.


\subsection{Entorno computacional}
\label{subsec:entorno_computacional}

La implementación se desarrolló en el entorno resumido en la
Tabla~\ref{tab:entorno_computacional}.

\begin{table}[htbp]
    \centering
    \caption{Entorno utilizado en los experimentos.}
    \label{tab:entorno_computacional}
    \begin{tabular}{p{0.30\textwidth}p{0.60\textwidth}}
        \hline
        \textbf{Componente} & \textbf{Configuración} \\
        \hline
        Sistema operativo & Microsoft Windows 10 \\
        Lenguaje & Python 3.12.10 \\
        Modelado matemático & PuLP \\
        Solver & CBC mediante la interfaz de PuLP \\
        Control de versiones & Git y GitHub \\
        Informe & \LaTeX, \texttt{biblatex} y Biber \\
        \hline
    \end{tabular}
\end{table}

La información específica de cada ejecución se almacena en

\begin{center}
\texttt{outputs/results/experiment\_environment.json}.
\end{center}

Este archivo registra las versiones disponibles, el modo de ejecución y
los datos de trazabilidad del experimento. Los tiempos observados se
consideran referencias operativas y pueden variar según el equipo, la
versión del solver y la carga del sistema.


\subsection{Procedimiento de reproducción}
\label{subsec:procedimiento_reproduccion}

Desde la raíz del repositorio, la reproducción mínima requiere ejecutar
la suite de pruebas y regenerar los artefactos experimentales:

\begin{verbatim}
python -m pytest -q --disable-warnings
python scripts/run_active_sbox_experiments.py --mode quick
\end{verbatim}

El modo \texttt{quick} utiliza los resultados permanentes de las búsquedas
más costosas y regenera:

\begin{itemize}
    \item la tabla consolidada de experimentos;
    \item los testigos almacenados;
    \item el resumen de resultados;
    \item la información del entorno;
    \item la figura comparativa de cotas.
\end{itemize}

Para repetir desde cero la búsqueda restringida de tres rondas se utiliza:

\begin{verbatim}
python scripts/run_active_sbox_experiments.py --mode full
\end{verbatim}

Este modo vuelve a evaluar las trayectorias de la familia

\[
2+2+c
\]

para $z=4$ y $z=8$.

El informe se compila desde la carpeta \texttt{report/} mediante:

\begin{verbatim}
latexmk -pdf -interaction=nonstopmode -halt-on-error main.tex
\end{verbatim}

Si la bibliografía se gestiona mediante Biber, \texttt{latexmk} ejecuta
automáticamente las etapas necesarias siempre que el archivo
\texttt{.bcf} se haya generado correctamente.


\subsection{Pruebas automatizadas}
\label{subsec:pruebas_automatizadas}

La suite global de la versión utilizada en el informe alcanzó

\[
245
\]

pruebas aprobadas.

Las pruebas cubren los siguientes aspectos:

\begin{itemize}
    \item transformaciones individuales de Keccak reducido;
    \item composición de rondas completas;
    \item restricciones XOR y productos binarios;
    \item correspondencia entre la implementación funcional y el modelo
    MILP;
    \item cálculo de diferencias y soportes activos;
    \item restricciones de actividad total y actividad por ronda;
    \item búsqueda restringida de trayectorias $2+2+c$;
    \item reconstrucción de testigos;
    \item generación y lectura de los artefactos experimentales.
\end{itemize}

La aprobación de las pruebas reduce el riesgo de errores de implementación,
pero no constituye una demostración formal de corrección completa del
software.

Durante la ejecución pueden aparecer advertencias de PuLP relacionadas con
interfaces marcadas para futura obsolescencia. Estas advertencias no
corresponden a fallos de factibilidad, optimalidad o reconstrucción de los
resultados.


\subsection{Reconstrucción de testigos}
\label{subsec:reconstruccion_testigos}

Cada testigo utilizado en los resultados contiene información suficiente
para recuperar los estados iniciales

\[
A_0^L
\qquad \text{y} \qquad
A_0^R.
\]

A partir de estos estados se ejecutan nuevamente las rondas mediante la
implementación funcional. Para cada ronda se calcula

\[
\Delta B_r
=
B_r^L \oplus B_r^R
\]

y se reconstruye el soporte activo

\[
S_r
=
\left\{
(y,k):
\bigvee_{x=0}^{4}
\Delta B_r[x,y,k]=1
\right\}.
\]

La actividad por ronda se obtiene como

\[
m_r=|S_r|.
\]

Un testigo se considera validado cuando se cumplen simultáneamente las
siguientes condiciones:

\begin{enumerate}
    \item los estados intermedios coinciden con las transformaciones de
    referencia;

    \item las diferencias almacenadas coinciden con los XOR reconstruidos;

    \item los soportes activos coinciden con las posiciones no nulas de
    $\Delta B_r$;

    \item la suma

    \[
    \sum_{r=0}^{R-1}m_r
    \]

    coincide con la actividad total reportada;

    \item los estados y soportes son consistentes con las restricciones del
    modelo MILP.
\end{enumerate}

Esta reconstrucción permite verificar independientemente los resultados
exactos de una y dos rondas y los testigos utilizados como cotas superiores
en tres rondas.


\subsection{Artefactos experimentales}
\label{subsec:artefactos_experimentales}

Los resultados se almacenan en archivos estructurados para evitar la
transcripción manual de valores. La
Tabla~\ref{tab:artefactos_experimentales} resume los artefactos principales.

\begin{table}[htbp]
    \centering
    \caption{Artefactos principales para reproducción y auditoría.}
    \label{tab:artefactos_experimentales}
    \begin{tabular}{p{0.42\textwidth}p{0.50\textwidth}}
        \hline
        \textbf{Artefacto} & \textbf{Contenido} \\
        \hline

        \texttt{active\_sbox\_results.csv}
        &
        Cotas, actividades y clasificación de los seis experimentos.
        \\[3pt]

        \texttt{active\_sbox\_witnesses.json}
        &
        Estados, diferencias y soportes necesarios para reconstruir los
        testigos.
        \\[3pt]

        \texttt{active\_sbox\_summary.md}
        &
        Resumen legible de los resultados principales.
        \\[3pt]

        \texttt{experiment\_environment.json}
        &
        Versiones, modo de ejecución y datos del entorno.
        \\[3pt]

        \texttt{active\_sbox\_bounds.png}
        &
        Comparación gráfica de las cotas inferiores y superiores.
        \\[3pt]

        \texttt{keccak\_milp\_active\_sboxes.ipynb}
        &
        Notebook ejecutable con tablas, figura y comprobaciones dirigidas.
        \\
        \hline
    \end{tabular}
\end{table}

Los archivos de resultados se encuentran dentro de
\texttt{outputs/results/}, la figura dentro de
\texttt{outputs/figures/} y el notebook dentro de
\texttt{notebooks/}.

La separación entre resultados numéricos, testigos y datos del entorno
facilita la auditoría y permite reutilizar los artefactos desde otros
scripts.


\subsection{Validación del notebook}
\label{subsec:validacion_notebook}

El notebook final fue ejecutado completamente antes de generar la versión
utilizada en el informe. La revisión comprobó que:

\begin{itemize}
    \item no existieran errores almacenados;
    \item todas las celdas de código hubieran sido ejecutadas;
    \item estuvieran incluidos los resultados para $z=4$ y $z=8$;
    \item aparecieran los testigos de tres rondas;
    \item las tablas y la figura coincidieran con los archivos generados
    por el ejecutor experimental.
\end{itemize}

El notebook funciona como una interfaz de inspección de los resultados,
pero los valores principales provienen de los archivos estructurados
generados por los scripts del proyecto.


\subsection{Trazabilidad mediante Git}
\label{subsec:trazabilidad_git}

La versión experimental utilizada en este informe se identifica mediante:

\begin{itemize}
    \item repositorio:

    \begin{center}
    \url{https://github.com/hsancheza90-art/keccak-milp-criptografia}
    \end{center}

    \item etiqueta:

    \begin{center}
    \texttt{keccak-milp-experiments-v1}
    \end{center}

    \item commit:

    \begin{center}
    \texttt{0b089aa}
    \end{center}
\end{itemize}

La etiqueta identifica un estado estable del proyecto y permite separar los
resultados utilizados en el informe de cambios posteriores realizados en la
rama de desarrollo.


\subsection{Alcance de la reproducibilidad}
\label{subsec:alcance_reproducibilidad}

El modo rápido permite regenerar los resultados permanentes y verificar la
cadena completa de archivos sin repetir las enumeraciones de mayor costo.
El modo completo vuelve a ejecutar la búsqueda restringida de tres rondas.

En ambos casos pueden verificarse:

\begin{itemize}
    \item los mínimos globales de una ronda;
    \item los mínimos globales de dos rondas;
    \item las cotas inferiores de tres rondas;
    \item los testigos $2+2+9$ y $2+2+10$;
    \item la consistencia entre los archivos, el notebook y el informe.
\end{itemize}

La reproducibilidad no implica que los tiempos de ejecución sean idénticos
en todos los equipos. Tampoco convierte los testigos de tres rondas en
óptimos globales. Su propósito es garantizar que las afirmaciones puedan
ser reconstruidas y contrastadas a partir del código y de los artefactos
publicados.